% =============================================================================
% APPENDIX BN: REF-SUPERKEY: Database Superkey Specification
% =============================================================================
% Category: REF (Reference)
% Module: BCM2_00_META
% Version: 1.0 (January 2026)
% Status: Final
% Language: English
%
% This appendix defines the unified superkey system for all EBF databases,
% ensuring referential integrity and consistent cross-database linking.
%
% =============================================================================

% =============================================================================
% CROSS-REFERENCE STRUCTURE
% =============================================================================
%
% CHAPTER LINKAGE (Required):
% - Primary Chapter: N/A (Meta-Infrastructure)
% - Secondary Chapters: All chapters that reference data registries
%
% DEPENDENCIES (what this appendix REQUIRES):
% - None (foundational specification)
%
% DEPENDENTS (what REQUIRES this appendix):
% - All data registries (model-registry.yaml, session-registry.yaml, etc.)
% - Appendix BBB (CORE-WHERE): Parameter references
% - Appendix G (Glossary): Symbol definitions
%
% =============================================================================

\chapter{REF-SUPERKEY: Database Superkey Specification}
\label{app:superkey}


% -----------------------------------------------------------------------------
% HEADER BLOCK
% -----------------------------------------------------------------------------

\begin{tcolorbox}[colback=blue!5!white, colframe=blue!75!black,
    title=Appendix Metadata]

\textbf{Appendix:} BN REF-SUPERKEY

\textbf{Category:} REF (Reference)

\textbf{Version:} 1.0 (January 2026)

\textbf{Status:} Final

\textbf{Language:} English

\textbf{Purpose:} Defines the unified superkey system for all EBF databases, ensuring referential integrity and consistent cross-database linking.

\end{tcolorbox}


% -----------------------------------------------------------------------------
% CHAPTER LINKAGE BOX
% -----------------------------------------------------------------------------

\begin{tcolorbox}[colback=purple!5!white, colframe=purple!75!black,
    title=Chapter Linkage]

\textbf{Primary Chapter:} Meta-Infrastructure (no dedicated chapter)

\textbf{Secondary Chapters:}
\begin{itemize}[nosep]
\item All chapters --- reference entities via superkeys
\item CLAUDE.md --- operational implementation
\end{itemize}

\textbf{Reading Guide:}
\begin{itemize}[nosep]
\item For conceptual overview: Read Section~\ref{sec:superkey-fundamental} first
\item For implementation: See Section~\ref{sec:superkey-specification}
\item For migration: See Section~\ref{sec:superkey-migration}
\end{itemize}

\end{tcolorbox}


% -----------------------------------------------------------------------------
% ABSTRACT
% -----------------------------------------------------------------------------

\begin{abstract}
This appendix specifies the unified superkey system for all EBF databases.
A superkey is a unique identifier that enables cross-database referencing
and ensures referential integrity. The specification covers 13 entity types
across 5 hierarchical levels: Organisation, Knowledge, Output, Action, and Learning.
Key design decisions include: hierarchical keys for customer-bound entities,
flat keys for global knowledge entities, version management via attributes
rather than key modification, and PAP- prefix for papers to prevent LLM hallucination.
\end{abstract}


% -----------------------------------------------------------------------------
% QUICK REFERENCE BOX
% -----------------------------------------------------------------------------

\begin{tcolorbox}[colback=blue!5!white, colframe=blue!75!black,
    title=Quick Reference: Superkey Specification]

\textbf{Core Question:} How do we uniquely identify and link entities across EBF databases?

\textbf{Key Pattern:}
\begin{equation}
\text{SUPERKEY} = \text{PREFIX} + \text{[-HIERARCHY]} + \text{-SEQUENCE}
\end{equation}

\textbf{Entity Types (13):}
\begin{itemize}[nosep]
\item \textbf{Organisation:} CUS (Customer), PRJ (Project)
\item \textbf{Knowledge:} EBF-S (Session), MOD (Model), MS (Theory), PAR (Parameter), CAS (Case), CON (Concept), FRM (Formula), PAP (Paper)
\item \textbf{Output:} OUT (Output)
\item \textbf{Action:} INT (Intervention)
\item \textbf{Learning:} INO-L (Innosuisse Learning)
\end{itemize}

\textbf{Cross-References:} Appendix G (Glossary), Appendix BBB (CORE-WHERE)
\end{tcolorbox}


% -----------------------------------------------------------------------------
% SCOPE BOX
% -----------------------------------------------------------------------------

\begin{tcolorbox}[
    colback=orange!5!white,
    colframe=orange!75!black,
    title={\textbf{Appendix Scope: Ziel / In-Scope / Out-of-Scope / Constraints}},
    fonttitle=\bfseries
]

\textbf{Ziel (Objective):}
\begin{itemize}[nosep]
    \item Define a unified, consistent superkey system for all EBF databases
\end{itemize}

\textbf{In-Scope:}
\begin{itemize}[nosep]
    \item Superkey format specification for all 13 entity types
    \item Naming conventions and syntax rules
    \item Cross-reference patterns between entities
    \item Version management strategy
    \item Migration plan from legacy identifiers
\end{itemize}

\textbf{Out-of-Scope (delegiert an andere Appendices):}
\begin{itemize}[nosep]
    \item Database schema details $\rightarrow$ Individual registry files
    \item Parameter values $\rightarrow$ Appendix BBB (CORE-WHERE)
    \item Symbol definitions $\rightarrow$ Appendix G (Glossary)
\end{itemize}

\textbf{Constraints (Anwendungsgrenzen):}
\begin{itemize}[nosep]
    \item All new entities MUST use this specification
    \item Legacy identifiers require migration (see Section~\ref{sec:superkey-migration})
    \item Superkeys are immutable once assigned (no renaming)
\end{itemize}

\textbf{Lieferobjekte (Deliverables):}
\begin{itemize}[nosep]
    \item Complete superkey format table (Table~\ref{tab:superkey-formats})
    \item Axiom system (SK-1 to SK-9)
    \item Migration mapping (Table~\ref{tab:superkey-migration})
    \item Validation rules for automated checking
\end{itemize}

\end{tcolorbox}


%%%%%%%%%%%%%%%%%%%%%%%%%%%%%%%%%%%%%%%%%%%%%%%%%%%%%%%%%%%%%%%%%%%%%%%%%%%%%%%
\section{The Fundamental Question}
\label{sec:superkey-fundamental}
%%%%%%%%%%%%%%%%%%%%%%%%%%%%%%%%%%%%%%%%%%%%%%%%%%%%%%%%%%%%%%%%%%%%%%%%%%%%%%%

\subsection{Why This Matters}

The EBF framework operates across multiple interconnected databases:
\begin{itemize}
    \item \texttt{model-registry.yaml} --- Behavioral models
    \item \texttt{model-building-session.yaml} --- Analysis sessions
    \item \texttt{parameter-registry.yaml} --- Behavioral parameters
    \item \texttt{intervention-registry.yaml} --- Projects and interventions
    \item \texttt{case-registry.yaml} --- Precedent cases
    \item \texttt{output-registry.yaml} --- Generated deliverables
    \item \texttt{concept-registry.yaml} --- EIP-validated concepts
    \item \texttt{formula-registry.yaml} --- Mathematical formulas
    \item \texttt{theory-catalog.yaml} --- Scientific theories
    \item \texttt{innosuisse-learning-database.yaml} --- Project-specific learnings
    \item Customer databases (\texttt{data/customers/})
\end{itemize}

Without a unified identifier system, cross-referencing becomes error-prone
and data integrity cannot be guaranteed.

\subsection{The Core Problem}

\begin{tcolorbox}[colback=red!5!white, colframe=red!75!black,
    title=The Fundamental Question]
\textbf{How do we uniquely identify every entity in the EBF ecosystem
and maintain referential integrity across all databases?}
\end{tcolorbox}


%%%%%%%%%%%%%%%%%%%%%%%%%%%%%%%%%%%%%%%%%%%%%%%%%%%%%%%%%%%%%%%%%%%%%%%%%%%%%%%
\section{Core Theory}
\label{sec:superkey-theory}
%%%%%%%%%%%%%%%%%%%%%%%%%%%%%%%%%%%%%%%%%%%%%%%%%%%%%%%%%%%%%%%%%%%%%%%%%%%%%%%

This section defines the theoretical foundation of the superkey system.

%%%%%%%%%%%%%%%%%%%%%%%%%%%%%%%%%%%%%%%%%%%%%%%%%%%%%%%%%%%%%%%%%%%%%%%%%%%%%%%
\section{Superkey Specification}
\label{sec:superkey-specification}
%%%%%%%%%%%%%%%%%%%%%%%%%%%%%%%%%%%%%%%%%%%%%%%%%%%%%%%%%%%%%%%%%%%%%%%%%%%%%%%

\subsection{Design Principles}

The superkey system follows five design principles:

\begin{enumerate}
    \item \textbf{Uniqueness:} Every superkey identifies exactly one entity
    \item \textbf{Stability:} Superkeys never change after assignment
    \item \textbf{Readability:} Superkeys are human-readable and meaningful
    \item \textbf{Hierarchy:} Customer-bound entities embed customer context
    \item \textbf{Simplicity:} Prefer short keys with attributes over long keys
\end{enumerate}

\subsection{Entity Hierarchy}

\begin{figure}[htbp]
\centering
\begin{verbatim}
┌─────────────────────────────────────────────────────────────────────────┐
│  LEVEL 1: ORGANISATION                                                  │
│  ─────────────────────                                                  │
│  CUS-{CODE}                     Customer    CUS-ALPLA, CUS-UBS          │
│      │                                                                  │
│      └──► CUS-{CODE}-PRJ-{NNN}  Project     CUS-ALPLA-PRJ-001           │
│                                                                         │
│  LEVEL 2: KNOWLEDGE (global, flat)                                      │
│  ─────────────────────────────────                                      │
│  EBF-S-{YYYY-MM-DD}-{DOM}-{NN}  Session     EBF-S-2026-01-25-REL-01     │
│  MOD-{NNN|SEM}                  Model       MOD-001, MOD-PSF            │
│  MS-{CAT}-{NNN}                 Theory      MS-RD-001 (historical)      │
│  PAR-{TYP}-{NNN}                Parameter   PAR-BEH-001                 │
│  CAS-{NNN}                      Case        CAS-001                     │
│  CON-{YYYY}-{NNN}               Concept     CON-2026-001                │
│  FRM-{NNN}                      Formula     FRM-001                     │
│  PAP-{BibKey}                   Paper       PAP-kahneman1979prospect    │
│                                                                         │
│  LEVEL 3: OUTPUT (session-bound via attribute)                          │
│  ─────────────────────────────────────────────                          │
│  OUT-{NNN}                      Output      OUT-001                     │
│      → session_ref: EBF-S-...                                           │
│                                                                         │
│  LEVEL 4: ACTION (project-bound)                                        │
│  ───────────────────────────────                                        │
│  CUS-{CODE}-PRJ-{NNN}-INT-{NN}  Intervention CUS-ALPLA-PRJ-001-INT-01   │
│                                                                         │
│  LEVEL 5: LEARNING (project-specific knowledge capture)                 │
│  ──────────────────────────────────────────────────────                 │
│  INO-L-{YYYY-MM-DD}-{CAT}-{NNN} Learning     INO-L-2026-01-27-DOC-001   │
│                                                                         │
└─────────────────────────────────────────────────────────────────────────┘
\end{verbatim}
\caption{Superkey Hierarchy}
\label{fig:superkey-hierarchy}
\end{figure}

\subsection{Complete Format Specification}

\begin{table}[htbp]
\centering
\caption{Superkey Formats by Entity Type}
\label{tab:superkey-formats}
\renewcommand{\arraystretch}{1.4}
\begin{tabular}{@{}llllp{3.5cm}@{}}
\toprule
\textbf{Entity} & \textbf{Prefix} & \textbf{Format} & \textbf{Example} & \textbf{Notes} \\
\midrule
\multicolumn{5}{l}{\textit{Organisation Level}} \\
Customer & CUS & CUS-\{CODE\} & CUS-ALPLA & CODE: 2-6 chars, natural \\
Project & PRJ & CUS-\{CODE\}-PRJ-\{NNN\} & CUS-ALPLA-PRJ-001 & Hierarchical under customer \\
\midrule
\multicolumn{5}{l}{\textit{Knowledge Level}} \\
Session & EBF-S & EBF-S-\{DATE\}-\{DOM\}-\{NN\} & EBF-S-2026-01-25-REL-01 & Date: YYYY-MM-DD \\
Model & MOD & MOD-\{NNN$|$SEM\} & MOD-001, MOD-PSF & Semantic IDs allowed \\
Theory & MS & MS-\{CAT\}-\{NNN\} & MS-RD-001 & Historical exception \\
Parameter & PAR & PAR-\{TYP\}-\{NNN\} & PAR-BEH-001 & TYP: BEH, CTX, INT, CMP \\
Case & CAS & CAS-\{NNN\} & CAS-001 & 3-digit sequence \\
Concept & CON & CON-\{YYYY\}-\{NNN\} & CON-2026-001 & Year of creation \\
Formula & FRM & FRM-\{NNN\} & FRM-001 & 3-digit sequence \\
Paper & PAP & PAP-\{BibKey\} & PAP-kahneman1979prospect & Prevents hallucination \\
\midrule
\multicolumn{5}{l}{\textit{Output Level}} \\
Output & OUT & OUT-\{NNN\} & OUT-001 & + session\_ref attribute \\
\midrule
\multicolumn{5}{l}{\textit{Action Level}} \\
Intervention & INT & CUS-\{CODE\}-PRJ-\{NNN\}-INT-\{NN\} & CUS-ALPLA-PRJ-001-INT-01 & Hierarchical under project \\
\midrule
\multicolumn{5}{l}{\textit{Learning Level}} \\
Learning & INO-L & INO-L-\{DATE\}-\{CAT\}-\{NNN\} & INO-L-2026-01-27-DOC-001 & Date: YYYY-MM-DD, CAT: category \\
\bottomrule
\end{tabular}
\end{table}

\subsection{Domain Codes}

Session superkeys include a domain code:

\begin{table}[htbp]
\centering
\caption{Domain Codes for Sessions}
\label{tab:domain-codes}
\begin{tabular}{@{}ll@{}}
\toprule
\textbf{Code} & \textbf{Domain} \\
\midrule
REL & Religion \\
FIN & Finance \\
HLT & Health \\
ENV & Environment \\
POL & Policy \\
ORG & Organisation \\
EDU & Education \\
COG & Cognitive \\
OTH & Other \\
\bottomrule
\end{tabular}
\end{table}

\subsection{Parameter Type Codes}

Parameter superkeys include a type code:

\begin{table}[htbp]
\centering
\caption{Parameter Type Codes}
\label{tab:param-codes}
\begin{tabular}{@{}ll@{}}
\toprule
\textbf{Code} & \textbf{Type} \\
\midrule
BEH & Behavioral (λ, β, γ, etc.) \\
CTX & Contextual (τ, trust, etc.) \\
INT & Intervention effect sizes (E\_i) \\
CMP & Complementarity (γ\_ij) \\
\bottomrule
\end{tabular}
\end{table}

\subsection{Learning Category Codes}

Learning superkeys include a category code:

\begin{table}[htbp]
\centering
\caption{Learning Category Codes}
\label{tab:learning-codes}
\begin{tabular}{@{}ll@{}}
\toprule
\textbf{Code} & \textbf{Category} \\
\midrule
DOC & Dokument-Erstellung und -Adaption \\
PROC & Prozess und Workflow \\
TECH & Technische Aspekte \\
COMM & Kommunikation und Feedback \\
QA & Qualitätssicherung \\
MEET & Meeting-Verarbeitung \\
\bottomrule
\end{tabular}
\end{table}


%%%%%%%%%%%%%%%%%%%%%%%%%%%%%%%%%%%%%%%%%%%%%%%%%%%%%%%%%%%%%%%%%%%%%%%%%%%%%%%
\section{Axioms and Formal Definitions}
\label{sec:superkey-axioms}
%%%%%%%%%%%%%%%%%%%%%%%%%%%%%%%%%%%%%%%%%%%%%%%%%%%%%%%%%%%%%%%%%%%%%%%%%%%%%%%

\begin{tcolorbox}[colback=gray!5!white, colframe=gray!75!black,
    title=Axiom SK-1: Uniqueness]
\textbf{Statement:} $\forall e_1, e_2 \in \mathcal{E}: \text{SK}(e_1) = \text{SK}(e_2) \Rightarrow e_1 = e_2$

\textbf{Interpretation:} Every superkey identifies exactly one entity.

\textbf{Implication:} Duplicate superkeys are forbidden; validation must reject them.
\end{tcolorbox}

\begin{tcolorbox}[colback=gray!5!white, colframe=gray!75!black,
    title=Axiom SK-2: Immutability]
\textbf{Statement:} $\forall e \in \mathcal{E}, \forall t_1 < t_2: \text{SK}(e, t_1) = \text{SK}(e, t_2)$

\textbf{Interpretation:} Superkeys never change after assignment.

\textbf{Implication:} Renaming requires deprecation + new entity creation.
\end{tcolorbox}

\begin{tcolorbox}[colback=gray!5!white, colframe=gray!75!black,
    title=Axiom SK-3: Prefix Determinism]
\textbf{Statement:} $\text{Type}(e) = f(\text{Prefix}(\text{SK}(e)))$

\textbf{Interpretation:} The prefix uniquely determines the entity type.

\textbf{Implication:} CUS- is always a customer, MOD- is always a model.
\end{tcolorbox}

\begin{tcolorbox}[colback=gray!5!white, colframe=gray!75!black,
    title=Axiom SK-4: Hierarchical Embedding]
\textbf{Statement:} $\text{Parent}(e) = \text{Extract}(\text{SK}(e))$ for hierarchical entities

\textbf{Interpretation:} Child entities embed their parent's superkey.

\textbf{Implication:} CUS-ALPLA-PRJ-001 → parent is CUS-ALPLA.
\end{tcolorbox}

\begin{tcolorbox}[colback=gray!5!white, colframe=gray!75!black,
    title=Axiom SK-5: Separator Consistency]
\textbf{Statement:} $\text{Separator} = \text{``-''} \quad \forall \text{SK}$

\textbf{Interpretation:} All superkeys use hyphen as the only separator.

\textbf{Implication:} No underscores, dots, or other separators allowed.
\end{tcolorbox}

\begin{tcolorbox}[colback=gray!5!white, colframe=gray!75!black,
    title=Axiom SK-6: Version Externality]
\textbf{Statement:} $\text{Version}(e) \notin \text{SK}(e)$

\textbf{Interpretation:} Version information is stored as an attribute, not in the superkey.

\textbf{Implication:} MOD-PSF has version: 2.0 as attribute, not MOD-PSF-2.
\end{tcolorbox}

\begin{tcolorbox}[colback=gray!5!white, colframe=gray!75!black,
    title=Axiom SK-7: Reference Validity]
\textbf{Statement:} $\forall \text{ref} \in \text{DB}: \exists e \in \mathcal{E}: \text{SK}(e) = \text{ref}$

\textbf{Interpretation:} Every reference must point to an existing entity.

\textbf{Implication:} Dangling references are validation errors.
\end{tcolorbox}

\begin{tcolorbox}[colback=gray!5!white, colframe=gray!75!black,
    title=Axiom SK-8: Historical Exception]
\textbf{Statement:} $\text{Prefix}(\text{Theory}) = \text{``MS-''} \neq \text{``THE-''}$

\textbf{Interpretation:} Theories retain the historical MS- prefix.

\textbf{Implication:} No migration required for theory-catalog.yaml.
\end{tcolorbox}

\begin{tcolorbox}[colback=gray!5!white, colframe=gray!75!black,
    title=Axiom SK-9: Paper Exception]
\textbf{Statement:} $\forall p \in \mathcal{P}: \text{SK}(p) = \text{``PAP-''} + \text{BibKey}(p)$

\textbf{Interpretation:} Papers use PAP- prefix followed by the BibTeX key.

\textbf{Rationale:} Prevents LLM hallucination by making paper references unambiguous.
Without PAP- prefix, an LLM might invent plausible-looking BibTeX keys (e.g.,
``smith2020theory'') that don't exist in the bibliography.

\textbf{Format:} PAP-\{author\}\{year\}\{keyword\}

\textbf{Examples:}
\begin{itemize}[nosep]
\item kahneman1979prospect $\rightarrow$ PAP-kahneman1979prospect
\item fehr1999theory $\rightarrow$ PAP-fehr1999theory
\item thaler2008nudge $\rightarrow$ PAP-thaler2008nudge
\end{itemize}

\textbf{Implication:} All BibTeX entries in bcm\_master.bib must be migrated to PAP- prefix.
\end{tcolorbox}

\begin{tcolorbox}[colback=gray!5!white, colframe=gray!75!black,
    title=Axiom SK-10: Learning Database Keys]
\textbf{Statement:} $\forall l \in \mathcal{L}: \text{SK}(l) = \text{``INO-L-''} + \text{DATE}(l) + \text{``-''} + \text{CAT}(l) + \text{``-''} + \text{SEQ}(l)$

\textbf{Interpretation:} Learning entries use INO-L prefix followed by date, category, and sequence.

\textbf{Rationale:} Project-specific learning databases capture knowledge from practical work
(document creation, meetings, process improvements). The date-first format enables chronological
ordering while the category enables filtering by learning type.

\textbf{Format:} INO-L-\{YYYY-MM-DD\}-\{CAT\}-\{NNN\}

\textbf{Categories:}
\begin{itemize}[nosep]
\item DOC --- Dokument-Erstellung und -Adaption
\item PROC --- Prozess und Workflow
\item TECH --- Technische Aspekte
\item COMM --- Kommunikation und Feedback
\item QA --- Qualitätssicherung
\item MEET --- Meeting-Verarbeitung
\end{itemize}

\textbf{Examples:}
\begin{itemize}[nosep]
\item INO-L-2026-01-27-DOC-001 (Versionsnummern-Konsistenz)
\item INO-L-2026-01-27-PROC-001 (Ernst Fehr Feedback-Muster)
\item INO-L-2026-01-27-MEET-001 (Meeting-Transkript-Analyse)
\end{itemize}

\textbf{Implication:} All entries in innosuisse-learning-database.yaml must use this format.
\end{tcolorbox}


%%%%%%%%%%%%%%%%%%%%%%%%%%%%%%%%%%%%%%%%%%%%%%%%%%%%%%%%%%%%%%%%%%%%%%%%%%%%%%%
\section{Key Results}
\label{sec:superkey-results}
%%%%%%%%%%%%%%%%%%%%%%%%%%%%%%%%%%%%%%%%%%%%%%%%%%%%%%%%%%%%%%%%%%%%%%%%%%%%%%%

The superkey specification establishes a unified identification system across all EBF databases:

\begin{enumerate}
    \item \textbf{Unified Entity Identification:} 11 entity types are uniquely identified through prefix-based superkeys, enabling unambiguous cross-database referencing.

    \item \textbf{Hierarchical Organisation:} Customer-bound entities (Projects, Interventions) embed their parent key, creating natural hierarchies without separate lookup tables.

    \item \textbf{Knowledge Graph Integration:} Flat keys for knowledge entities (Models, Theories, Parameters) enable flexible many-to-many relationships across the EBF knowledge base.

    \item \textbf{Session Traceability:} Every analytical output can be traced to its originating session via the \texttt{session\_ref} attribute, supporting full audit trails.

    \item \textbf{Version Independence:} Semantic versioning as attributes (not in keys) ensures referential stability while supporting entity evolution.

    \item \textbf{Backward Compatibility:} Historical formats (MS- for theories) are preserved, requiring no migration of existing data.
\end{enumerate}

\paragraph{Quantitative Summary.}
\begin{itemize}
    \item 5 hierarchical levels (Organisation, Knowledge, Output, Action, Learning)
    \item 13 entity types with unique prefixes (including PAP- for papers, INO-L for learnings)
    \item 10 formal axioms guaranteeing consistency (SK-1 to SK-10)
    \item 10 cross-reference patterns defined
\end{itemize}


%%%%%%%%%%%%%%%%%%%%%%%%%%%%%%%%%%%%%%%%%%%%%%%%%%%%%%%%%%%%%%%%%%%%%%%%%%%%%%%
\section{Cross-Reference Patterns}
\label{sec:superkey-crossref}
%%%%%%%%%%%%%%%%%%%%%%%%%%%%%%%%%%%%%%%%%%%%%%%%%%%%%%%%%%%%%%%%%%%%%%%%%%%%%%%

\subsection{Reference Types}

\begin{table}[htbp]
\centering
\caption{Cross-Reference Patterns}
\label{tab:crossref-patterns}
\renewcommand{\arraystretch}{1.3}
\begin{tabular}{@{}lllp{4cm}@{}}
\toprule
\textbf{From} & \textbf{To} & \textbf{Attribute} & \textbf{Example} \\
\midrule
Session & Customer & customer\_ref & CUS-UBS \\
Session & Project & project\_ref & CUS-UBS-PRJ-001 \\
Session & Model & model\_ref & MOD-PSF \\
Output & Session & session\_ref & EBF-S-2026-01-25-REL-01 \\
Model & Theory & theory\_basis & [MS-RD-001, MS-SP-001] \\
Model & Session & created\_in\_session & EBF-S-2026-01-25-REL-01 \\
Parameter & Session & source\_session & EBF-S-2026-01-25-FIN-01 \\
Parameter & Model & source\_model & MOD-REF-001 \\
Intervention & Project & (embedded) & CUS-ALPLA-PRJ-001-INT-01 \\
Concept & Session & session & EBF-S-2026-01-19-OTH-01 \\
\bottomrule
\end{tabular}
\end{table}

\subsection{Cardinality}

\begin{itemize}
    \item Customer → Projects: 1:N
    \item Project → Interventions: 1:N
    \item Session → Outputs: 1:N
    \item Session → Models: N:M (sessions can use/create multiple models)
    \item Model → Theories: N:M
    \item Model → Parameters: 1:N (model defines parameters)
\end{itemize}


%%%%%%%%%%%%%%%%%%%%%%%%%%%%%%%%%%%%%%%%%%%%%%%%%%%%%%%%%%%%%%%%%%%%%%%%%%%%%%%
\section{Worked Example}
\label{sec:superkey-example}
%%%%%%%%%%%%%%%%%%%%%%%%%%%%%%%%%%%%%%%%%%%%%%%%%%%%%%%%%%%%%%%%%%%%%%%%%%%%%%%

\subsection{Complete Project Lifecycle}

\paragraph{Setup.} ALPLA engages FehrAdvice for a behavioral strategy project.

\paragraph{Step 1: Customer Registration.}
\begin{verbatim}
id: CUS-ALPLA
name: "ALPLA Werke Alwin Lehner GmbH & Co. KG"
code: ALPLA
created: 2026-01-15
\end{verbatim}

\paragraph{Step 2: Project Creation.}
\begin{verbatim}
id: CUS-ALPLA-PRJ-001
customer_ref: CUS-ALPLA
name: "Strategic Option Comparison"
start_date: 2026-01-15
\end{verbatim}

\paragraph{Step 3: Analysis Session.}
\begin{verbatim}
id: EBF-S-2026-01-20-ORG-01
customer_ref: CUS-ALPLA
project_ref: CUS-ALPLA-PRJ-001
question: "Which strategic option maximizes utility?"
mode: TIEF
\end{verbatim}

\paragraph{Step 4: Model Creation.}
\begin{verbatim}
id: MOD-SOCM
version: 1.0
name: "Strategy Option Comparison Model"
created_in_session: EBF-S-2026-01-20-ORG-01
theory_basis: [MS-IB-001, MS-SP-001]
\end{verbatim}

\paragraph{Step 5: Output Generation.}
\begin{verbatim}
id: OUT-001
session_ref: EBF-S-2026-01-20-ORG-01
type: F3
format: pdf
\end{verbatim}

\paragraph{Step 6: Intervention Design.}
\begin{verbatim}
id: CUS-ALPLA-PRJ-001-INT-01
type: nudge
subtype: default
target_10c: WHEN
\end{verbatim}

\paragraph{Result.} Complete traceability from customer to intervention:
\begin{verbatim}
CUS-ALPLA
    └── CUS-ALPLA-PRJ-001
            ├── EBF-S-2026-01-20-ORG-01
            │       ├── MOD-SOCM (created)
            │       └── OUT-001 (generated)
            └── CUS-ALPLA-PRJ-001-INT-01
\end{verbatim}


%%%%%%%%%%%%%%%%%%%%%%%%%%%%%%%%%%%%%%%%%%%%%%%%%%%%%%%%%%%%%%%%%%%%%%%%%%%%%%%
\section{Migration Plan}
\label{sec:superkey-migration}
%%%%%%%%%%%%%%%%%%%%%%%%%%%%%%%%%%%%%%%%%%%%%%%%%%%%%%%%%%%%%%%%%%%%%%%%%%%%%%%

\begin{table}[htbp]
\centering
\caption{Migration from Legacy to New Superkeys}
\label{tab:superkey-migration}
\renewcommand{\arraystretch}{1.3}
\begin{tabular}{@{}llll@{}}
\toprule
\textbf{Entity} & \textbf{Legacy} & \textbf{New} & \textbf{Action} \\
\midrule
Session & EBF-S-2026-01-25-REL-001 & EBF-S-2026-01-25-REL-01 & Shorten sequence \\
Model & EBF-MOD-001 & MOD-001 & Remove EBF- prefix \\
Model & PSF-2.0 & MOD-PSF & Standardize, version as attr \\
Model & ESL-META-001 & MOD-ESL & Standardize \\
Output & EBF-OUT-001 & OUT-001 & Remove EBF- prefix \\
Case & CASE-001 & CAS-001 & 3-char prefix \\
Concept & CONC-2026-001 & CON-2026-001 & 3-char prefix \\
Project & PRJ-001 & CUS-XXX-PRJ-001 & Add customer context \\
Theory & MS-RD-001 & MS-RD-001 & No change (SK-8) \\
Customer & (none) & CUS-ALPLA & \textbf{NEW} \\
Formula & (none) & FRM-001 & \textbf{NEW} \\
\bottomrule
\end{tabular}
\end{table}

\subsection{Migration Strategy}

\begin{enumerate}
    \item \textbf{Phase 1:} Create customer registry with CUS- keys
    \item \textbf{Phase 2:} Update project references to include customer
    \item \textbf{Phase 3:} Standardize model IDs (remove EBF-MOD- prefix)
    \item \textbf{Phase 4:} Shorten case/concept prefixes
    \item \textbf{Phase 5:} Add FRM- keys to formula registry
    \item \textbf{Phase 6:} Validation script to check all references
\end{enumerate}


\subsection{Paper Migration: Incremental Learning Algorithm}
\label{sec:paper-migration-algorithm}

The migration of 2,278 papers to PAP- superkeys uses an \textbf{incremental learning
algorithm} that prioritizes simple cases and accumulates learnings for complex ones.

\begin{tcolorbox}[colback=yellow!5!white, colframe=yellow!75!black,
    title=Paper Migration Algorithm (Incremental Learning)]

\textbf{Core Principle:} Start with papers having fewest references, validate completely,
learn patterns, apply learnings to complex cases.

\textbf{Pre-Migration Validation:}
\begin{enumerate}[nosep]
    \item Verify DOI exists (CrossRef API or manual lookup)
    \item Verify URL is accessible (doi.org resolver)
    \item PAP- prefix is only assigned to \textbf{validated} papers
\end{enumerate}

\textbf{Migration Steps (per paper):}
\begin{enumerate}[nosep]
    \item Update BibTeX entry: \texttt{@article\{key,} $\rightarrow$ \texttt{@article\{PAP-key,}
    \item Update all references in:
    \begin{itemize}[nosep]
        \item LaTeX files: \verb|\citep{key}| $\rightarrow$ \verb|\citep{PAP-key}|
        \item YAML files: \texttt{bib\_keys}, \texttt{use\_for}, \texttt{theory\_support}
    \end{itemize}
    \item \textbf{Validate:} \texttt{grep "key" | grep -v "PAP-key"} must return empty
    \item Log learnings to \texttt{data/paper\_migration\_learnings.yaml}
\end{enumerate}

\textbf{Priority Order:}
\begin{enumerate}[nosep]
    \item Papers with 0 references (BibTeX only) --- 239 papers
    \item Papers with 1 reference --- 1,321 papers
    \item Papers with 2--5 references --- 410 papers
    \item Papers with 6+ references --- 308 papers
\end{enumerate}

\textbf{Learning Loop:}
\begin{itemize}[nosep]
    \item Each migration records: files changed, patterns found, issues encountered
    \item Patterns become rules for automated migration
    \item Edge cases inform manual review requirements
\end{itemize}

\end{tcolorbox}

\textbf{Tooling:}
\begin{itemize}
    \item \texttt{scripts/migrate\_paper\_validated.py --analyze} --- Show migration status
    \item \texttt{scripts/migrate\_paper\_validated.py --check KEY} --- Check single paper
    \item \texttt{scripts/migrate\_paper\_validated.py --migrate KEY --execute} --- Migrate paper
    \item \texttt{scripts/migrate\_paper\_validated.py --validate KEY} --- Verify no orphans
\end{itemize}

\textbf{Learnings Log:} \texttt{data/paper\_migration\_learnings.yaml}


%%%%%%%%%%%%%%%%%%%%%%%%%%%%%%%%%%%%%%%%%%%%%%%%%%%%%%%%%%%%%%%%%%%%%%%%%%%%%%%
\section{Implications for Practice}
\label{sec:superkey-practice}
%%%%%%%%%%%%%%%%%%%%%%%%%%%%%%%%%%%%%%%%%%%%%%%%%%%%%%%%%%%%%%%%%%%%%%%%%%%%%%%

\begin{tcolorbox}[colback=green!5!white, colframe=green!75!black,
    title=Practical Implications]
\begin{enumerate}
\item \textbf{Always use the registry:} Never invent superkeys ad-hoc; check availability first
\item \textbf{Customer context first:} For client work, create CUS- entry before anything else
\item \textbf{Version as attribute:} Never put version numbers in superkeys
\item \textbf{Validate references:} Run \texttt{validate\_referential\_integrity.py} before commits
\item \textbf{Document in session:} Every session should record which entities it created/used
\end{enumerate}
\end{tcolorbox}


%%%%%%%%%%%%%%%%%%%%%%%%%%%%%%%%%%%%%%%%%%%%%%%%%%%%%%%%%%%%%%%%%%%%%%%%%%%%%%%
\section{Summary}
\label{sec:superkey-summary}
%%%%%%%%%%%%%%%%%%%%%%%%%%%%%%%%%%%%%%%%%%%%%%%%%%%%%%%%%%%%%%%%%%%%%%%%%%%%%%%

\begin{tcolorbox}[colback=gray!5!white, colframe=gray!75!black,
    title=Appendix BN Summary: Superkey Specification]

\textbf{Core Contribution:}
A unified superkey system for all EBF databases ensuring referential integrity
and consistent cross-database linking.

\textbf{Key Pattern:}
\begin{equation}
\text{SUPERKEY} = \text{PREFIX} + \text{[-HIERARCHY]} + \text{-SEQUENCE}
\end{equation}

\textbf{Key Insights:}
\begin{itemize}[nosep]
\item 13 entity types across 5 hierarchical levels
\item Hierarchical keys for customer-bound entities (CUS-, PRJ-, INT-)
\item Flat keys for global knowledge entities (MOD-, PAR-, CAS-, PAP-)
\item Version management via attributes, not key modification
\item Historical exception: Theories retain MS- prefix (SK-8)
\item Paper exception: PAP- prefix prevents LLM hallucination (SK-9)
\item Learning exception: INO-L- prefix with date-category format (SK-10)
\end{itemize}

\textbf{Next Steps:} Create customer-registry.yaml and migrate existing keys
\end{tcolorbox}


%%%%%%%%%%%%%%%%%%%%%%%%%%%%%%%%%%%%%%%%%%%%%%%%%%%%%%%%%%%%%%%%%%%%%%%%%%%%%%%
\section{Glossary of Symbols}
\label{sec:superkey-glossary}
%%%%%%%%%%%%%%%%%%%%%%%%%%%%%%%%%%%%%%%%%%%%%%%%%%%%%%%%%%%%%%%%%%%%%%%%%%%%%%%

\begin{tcolorbox}[colback=blue!5!white, colframe=blue!75!black]
\textbf{Central Reference:} For complete symbol definitions, see
\textbf{Appendix G} (Glossary).

This section lists only symbols introduced in this appendix.
\end{tcolorbox}

\begin{table}[htbp]
\centering
\caption{Symbols Introduced in Appendix BN}
\label{tab:superkey-symbols}
\renewcommand{\arraystretch}{1.3}
\begin{tabular}{@{}clp{6cm}c@{}}
\toprule
\textbf{Symbol} & \textbf{Name} & \textbf{Definition} & \textbf{Also in G?} \\
\midrule
SK$(e)$ & Superkey function & Returns the superkey of entity $e$ & NEW \\
$\mathcal{E}$ & Entity set & Set of all EBF entities & NEW \\
PREFIX & Prefix component & Type-identifying prefix (CUS-, MOD-, etc.) & NEW \\
\bottomrule
\end{tabular}
\end{table}


%%%%%%%%%%%%%%%%%%%%%%%%%%%%%%%%%%%%%%%%%%%%%%%%%%%%%%%%%%%%%%%%%%%%%%%%%%%%%%%
\section{Open Issues and Research Agenda}
\label{sec:superkey-open}
%%%%%%%%%%%%%%%%%%%%%%%%%%%%%%%%%%%%%%%%%%%%%%%%%%%%%%%%%%%%%%%%%%%%%%%%%%%%%%%

\begin{table}[htbp]
\centering
\caption{Research Roadmap for Superkey System}
\label{tab:superkey-roadmap}
\renewcommand{\arraystretch}{1.3}
\begin{tabular}{@{}clcl@{}}
\toprule
\textbf{\#} & \textbf{Issue} & \textbf{Priority} & \textbf{Status} \\
\midrule
1 & Customer registry implementation & HIGH & PENDING \\
2 & Migration script for legacy keys & HIGH & PENDING \\
3 & Validation script enhancement & MEDIUM & PARTIAL \\
4 & UUID fallback for external systems & LOW & OPEN \\
\bottomrule
\end{tabular}
\end{table}


%%%%%%%%%%%%%%%%%%%%%%%%%%%%%%%%%%%%%%%%%%%%%%%%%%%%%%%%%%%%%%%%%%%%%%%%%%%%%%%
\section{References}
\label{sec:superkey-references}
%%%%%%%%%%%%%%%%%%%%%%%%%%%%%%%%%%%%%%%%%%%%%%%%%%%%%%%%%%%%%%%%%%%%%%%%%%%%%%%

\begin{tcolorbox}[colback=blue!5!white, colframe=blue!75!black]
\textbf{Master Bibliography:} For complete reference details, see
\texttt{bcm\_master.bib} in the repository.

This appendix is a methodological specification with no external references.
\end{tcolorbox}

\nocite{bcm_master}


%%%%%%%%%%%%%%%%%%%%%%%%%%%%%%%%%%%%%%%%%%%%%%%%%%%%%%%%%%%%%%%%%%%%%%%%%%%%%%%
% CROSS-REFERENCE FOOTER
%%%%%%%%%%%%%%%%%%%%%%%%%%%%%%%%%%%%%%%%%%%%%%%%%%%%%%%%%%%%%%%%%%%%%%%%%%%%%%%

\vspace{1em}
\hrule
\vspace{0.5em}

\noindent\textit{Cross-references:}
Appendix G (Glossary),
Appendix BBB (CORE-WHERE),
Appendix VC (METHOD-VALIDATE).

\noindent\textit{Master Bibliography:} \texttt{bcm\_master.bib}


% =============================================================================
% END OF APPENDIX BN
% =============================================================================
